% !TEX root = ../main.tex
\chapter{Data Cleaning and Preparation}
\label{ch:nettoyage}

\begin{intuition}
It is estimated that data scientists spend about 80\,\% of their time cleaning and preparing data. Poorly prepared data inevitably lead to flawed models: ``Garbage in, garbage out.''
\end{intuition}

% ============================================================
\section{Why Clean Data?}
\index{data cleaning}

Real-world data are rarely clean. Common issues include:
\begin{itemize}
    \item \textbf{Missing values}\index{missing values} (faulty sensors, unfilled fields);
    \item \textbf{Duplicates}\index{duplicates} (multiple entries, table merges);
    \item \textbf{Incorrect types}\index{data types} (numbers stored as text);
    \item \textbf{Outliers}\index{outliers} (data entry errors);
    \item \textbf{Textual inconsistencies} (capitalization, spaces, abbreviations).
\end{itemize}

% ============================================================
\section{Cleaning Pipeline}

\begin{center}
\begin{tikzpicture}[
    node distance=0.9cm,
    block/.style={rectangle, draw=teal!80, fill=teal!10, rounded corners, text width=4cm, align=center, minimum height=0.9cm, font=\small},
    arrow/.style={->, thick, >=stealth, color=gray!70}
]
\node[block] (raw) {Raw data};
\node[block, below=of raw] (dup) {Remove duplicates};
\node[block, below=of dup] (types) {Fix data types};
\node[block, below=of types] (miss) {Handle missing values};
\node[block, below=of miss] (out) {Handle outliers};
\node[block, below=of out] (text) {Text cleaning};
\node[block, below=of text] (feat) {Feature engineering};
\node[block, below=of feat] (enc) {Encoding and scaling};
\node[block, below=of enc] (clean) {Clean data};

\foreach \a/\b in {raw/dup, dup/types, types/miss, miss/out, out/text, text/feat, feat/enc, enc/clean}
    \draw[arrow] (\a) -- (\b);
\end{tikzpicture}
\end{center}

% ============================================================
\section{Case Study: The Titanic Dataset}
\index{Titanic}

\begin{codeblock}[title=\textbf{Python}]
\begin{minted}{python}
import pandas as pd
import numpy as np
import seaborn as sns

titanic = sns.load_dataset('titanic')
print(f"Initial dimensions: {titanic.shape}")
print(titanic.info())
\end{minted}
\end{codeblock}

\begin{codeoutput}
\begin{minted}{text}
Initial dimensions: (891, 15)
<class 'pandas.core.frame.DataFrame'>
RangeIndex: 891 entries, 0 to 890
Data columns (total 15 columns):
 #   Column       Non-Null Count  Dtype
 0   survived     891 non-null    int64
 1   pclass       891 non-null    int64
 2   sex          891 non-null    object
 3   age          714 non-null    float64
 ...
 7   embarked     889 non-null    object
 11  deck         203 non-null    category
\end{minted}
\end{codeoutput}

% ============================================================
\section{Missing Values}
\index{missing values}\index{MCAR}\index{MAR}\index{MNAR}

\subsection{Types of Missing Data}

\begin{definition}[Classification of Missing Values]
\begin{itemize}
    \item \textbf{MCAR} (\emph{Missing Completely At Random}): the missingness does not depend on any variable.
    \item \textbf{MAR} (\emph{Missing At Random}): the missingness depends on other observed variables.
    \item \textbf{MNAR} (\emph{Missing Not At Random}): the missingness depends on the value itself.
\end{itemize}
\end{definition}

\begin{example}
In the Titanic dataset, cabins (\py{deck}) are missing mainly for third-class passengers (MAR). Age could be MCAR or MAR depending on whether passengers without a full ticket had no recorded age.
\end{example}

\subsection{Detection and Treatment}

\begin{codeblock}[title=\textbf{Python}]
\begin{minted}{python}
# Detection
print(titanic.isnull().sum().sort_values(ascending=False).head(5))

# Strategy 1: drop the deck column (too many missing values)
titanic_clean = titanic.drop(columns=['deck'])

# Strategy 2: fill age with median by class
titanic_clean['age'] = titanic_clean.groupby('pclass')['age'] \
    .transform(lambda x: x.fillna(x.median()))

# Strategy 3: fill embarked with the mode
titanic_clean['embarked'].fillna(
    titanic_clean['embarked'].mode()[0], inplace=True)
titanic_clean['embark_town'].fillna(
    titanic_clean['embark_town'].mode()[0], inplace=True)

print(f"\nRemaining missing values: {titanic_clean.isnull().sum().sum()}")
\end{minted}
\end{codeblock}

\begin{codeoutput}
\begin{minted}{text}
deck           688
age            177
embarked         2
embark_town      2
survived         0

Remaining missing values: 0
\end{minted}
\end{codeoutput}

\begin{bestpractice}
Imputation by \emph{conditional} median (here by class) is preferable to using the global median because it preserves the structure of subgroups. For time series data, use \py{ffill()} or \py{bfill()}.
\end{bestpractice}

% ============================================================
\section{Duplicates}
\index{duplicates}

\begin{codeblock}[title=\textbf{Python}]
\begin{minted}{python}
# Detection
n_dup = titanic_clean.duplicated().sum()
print(f"Number of duplicates: {n_dup}")

# Inspect potential duplicates
if n_dup > 0:
    print(titanic_clean[titanic_clean.duplicated(keep=False)]
          .sort_values(by='name').head())
    titanic_clean = titanic_clean.drop_duplicates()
    print(f"After removal: {titanic_clean.shape}")
\end{minted}
\end{codeblock}

\begin{codeoutput}
\begin{minted}{text}
Number of duplicates: 0
\end{minted}
\end{codeoutput}

% ============================================================
\section{Type Conversion}
\index{type conversion}

\begin{codeblock}[title=\textbf{Python}]
\begin{minted}{python}
# Convert survived to boolean
titanic_clean['survived'] = titanic_clean['survived'].astype(bool)

# Convert pclass to ordered categorical
titanic_clean['pclass'] = pd.Categorical(
    titanic_clean['pclass'], categories=[1, 2, 3], ordered=True)

# Example of date conversion (with dummy data)
dates_str = pd.Series(['2023-01-15', '2023-02-20', '2023-03-10'])
dates = pd.to_datetime(dates_str)
print(dates.dt.month)
\end{minted}
\end{codeblock}

\begin{codeoutput}
\begin{minted}{text}
0    1
1    2
2    3
dtype: int32
\end{minted}
\end{codeoutput}

% ============================================================
\section{Handling Outliers}
\index{outliers}

\begin{codeblock}[title=\textbf{Python}]
\begin{minted}{python}
# IQR method for fare
Q1 = titanic_clean['fare'].quantile(0.25)
Q3 = titanic_clean['fare'].quantile(0.75)
IQR = Q3 - Q1

lower_bound = Q1 - 1.5 * IQR
upper_bound = Q3 + 1.5 * IQR

# Option 1: capping/winsorizing
titanic_clean['fare_capped'] = titanic_clean['fare'].clip(
    lower=lower_bound, upper=upper_bound)

print(f"Before: max = {titanic_clean['fare'].max():.2f}")
print(f"After capping: max = {titanic_clean['fare_capped'].max():.2f}")
print(f"Upper bound: {upper_bound:.2f}")
\end{minted}
\end{codeblock}

\begin{codeoutput}
\begin{minted}{text}
Before: max = 512.33
After capping: max = 65.63
Upper bound: 65.63
\end{minted}
\end{codeoutput}

\begin{warning}
Never remove outliers without careful thought. A fare of \$512 may be legitimate (luxury suite). Capping is often preferable to removal. Always document your choices.
\end{warning}

% ============================================================
\section{Text Cleaning}
\index{text cleaning}

\begin{codeblock}[title=\textbf{Python}]
\begin{minted}{python}
# Example with textual data
names = pd.Series(['  Alice ', 'BOB', 'charlie', ' Dave  '])
print("Before:", names.tolist())

names_clean = (names.str.strip()
                  .str.lower()
                  .str.capitalize())
print("After:", names_clean.tolist())

# Pattern replacement
cities = pd.Series(['New-York', 'new york', 'NEW YORK', 'N.Y.'])
cities_norm = cities.str.lower().str.replace(
    r'[^a-z\s]', '', regex=True).str.strip()
print("Normalized cities:", cities_norm.tolist())
\end{minted}
\end{codeblock}

\begin{codeoutput}
\begin{minted}{text}
Before: ['  Alice ', 'BOB', 'charlie', ' Dave  ']
After: ['Alice', 'Bob', 'Charlie', 'Dave']
Normalized cities: ['newyork', 'new york', 'new york', 'ny']
\end{minted}
\end{codeoutput}

% ============================================================
\section{Feature Engineering}
\index{feature engineering}

\begin{codeblock}[title=\textbf{Python}]
\begin{minted}{python}
# Create new variables from existing ones
titanic_clean['family_size'] = (titanic_clean['sibsp']
                                + titanic_clean['parch'] + 1)
titanic_clean['is_child'] = titanic_clean['age'] < 18
titanic_clean['fare_per_person'] = (titanic_clean['fare']
                                    / titanic_clean['family_size'])

# Discretize age into categories
titanic_clean['age_group'] = pd.cut(
    titanic_clean['age'],
    bins=[0, 12, 18, 35, 60, 100],
    labels=['Child', 'Teen', 'Young Adult', 'Adult', 'Senior'])

print(titanic_clean['age_group'].value_counts())
\end{minted}
\end{codeblock}

\begin{codeoutput}
\begin{minted}{text}
Young Adult    371
Adult          267
Child           89
Teen            79
Senior          85
Name: age_group, dtype: int64
\end{minted}
\end{codeoutput}

% ============================================================
\section{Encoding Categorical Variables}
\index{encoding}\index{one-hot encoding}\index{label encoding}

\begin{codeblock}[title=\textbf{Python}]
\begin{minted}{python}
# One-hot encoding with pandas
embarked_dummies = pd.get_dummies(
    titanic_clean['embarked'], prefix='embarked', dtype=int)
print(embarked_dummies.head())

# Label encoding with sklearn
from sklearn.preprocessing import LabelEncoder

le = LabelEncoder()
titanic_clean['sex_encoded'] = le.fit_transform(
    titanic_clean['sex'])
print("\nMapping:", dict(zip(le.classes_,
                              le.transform(le.classes_))))
\end{minted}
\end{codeblock}

\begin{codeoutput}
\begin{minted}{text}
   embarked_C  embarked_Q  embarked_S
0           0           0           1
1           1           0           0
2           0           0           1
3           0           0           1
4           0           0           1

Mapping: {'female': 0, 'male': 1}
\end{minted}
\end{codeoutput}

\begin{warning}
\textbf{Label encoding}\index{label encoding} imposes an artificial order on categories. It is only suitable for ordinal variables or tree-based algorithms. For linear models, prefer \textbf{one-hot encoding}\index{one-hot encoding}.
\end{warning}

% ============================================================
\section{Feature Scaling}
\index{standardization}\index{normalization}\index{StandardScaler}\index{MinMaxScaler}

\begin{definition}[Standardization and Normalization]
\begin{itemize}
    \item \textbf{Standardization} (StandardScaler): $z = \dfrac{x - \mu}{\sigma}$, resulting in mean 0 and standard deviation 1.
    \item \textbf{Normalization} (MinMaxScaler): $x' = \dfrac{x - x_{\min}}{x_{\max} - x_{\min}}$, resulting in values in $[0, 1]$.
\end{itemize}
\end{definition}

\begin{codeblock}[title=\textbf{Python}]
\begin{minted}{python}
from sklearn.preprocessing import StandardScaler, MinMaxScaler

cols = ['age', 'fare']
data_num = titanic_clean[cols].copy()

# Standardization
scaler_std = StandardScaler()
data_std = pd.DataFrame(scaler_std.fit_transform(data_num),
                        columns=[c + '_std' for c in cols])

# Normalization
scaler_mm = MinMaxScaler()
data_norm = pd.DataFrame(scaler_mm.fit_transform(data_num),
                         columns=[c + '_norm' for c in cols])

print("Standardized:")
print(data_std.describe().loc[['mean', 'std']].round(2))
print("\nNormalized:")
print(data_norm.describe().loc[['min', 'max']].round(2))
\end{minted}
\end{codeblock}

\begin{codeoutput}
\begin{minted}{text}
Standardized:
      age_std  fare_std
mean     0.00      0.00
std      1.00      1.00

Normalized:
      age_norm  fare_norm
min       0.00       0.00
max       1.00       1.00
\end{minted}
\end{codeoutput}

\begin{bestpractice}
Apply scaling \emph{after} the train/test split to avoid information leakage (\emph{data leakage}\index{data leakage}). Use \py{fit\_transform()} on the training set and \py{transform()} on the test set.
\end{bestpractice}

% ============================================================
\section{Exercises}

\begin{exercise}[$\star$]
Load the \py{titanic} dataset. Display the number and percentage of missing values per column. Which columns have more than 50\,\% missing values?
\end{exercise}

\begin{exercise}[$\star$]
Remove duplicates from the \py{titanic} dataset (if any). Then convert the \py{sex} column to a categorical variable with \py{pd.Categorical()}.
\end{exercise}

\begin{exercise}[$\star\star$]
Impute the missing values of \py{age} using the median conditioned on \py{sex} and \py{pclass} (6 groups). Compare the age distribution before and after imputation with overlapping histograms.
\end{exercise}

\begin{exercise}[$\star\star$]
Create the following variables: \py{title} (extracted from the name using a regular expression), \py{is\_alone} (true if \py{family\_size} = 1), \py{fare\_bin} (fare quartiles with \py{pd.qcut}). Compute the survival rate for each of these new variables.
\end{exercise}

\begin{exercise}[$\star\star\star$]
Build a complete cleaning pipeline for the \py{titanic} dataset: (a) drop \py{deck}, (b) impute \py{age} and \py{embarked}, (c) create \py{family\_size} and \py{is\_child}, (d) one-hot encode \py{sex}, \py{embarked}, \py{class}, (e) standardize numerical variables. The final result should be a fully numerical DataFrame with no missing values.
\end{exercise}

\begin{exercise}[$\star\star\star$]
Load the \py{penguins} dataset (\py{sns.load\_dataset('penguins')}). This dataset contains missing values. Apply three different strategies (deletion, median imputation, median imputation conditioned by species). Compare the descriptive statistics obtained with each strategy and discuss the best approach.
\end{exercise}

% ============================================================
\begin{keyformulas}
\textbf{Chapter Summary -- Data Cleaning and Preparation}
\begin{itemize}
    \item \textbf{Missing values}: \py{isnull().sum()}, MCAR/MAR/MNAR types, \py{fillna()}, \py{dropna()}.
    \item \textbf{Duplicates}: \py{duplicated()}, \py{drop\_duplicates()}.
    \item \textbf{Types}: \py{astype()}, \py{pd.to\_datetime()}, \py{pd.Categorical()}.
    \item \textbf{Outliers}: IQR method, capping with \py{clip()}.
    \item \textbf{Text}: \py{str.strip()}, \py{str.lower()}, \py{str.replace()}.
    \item \textbf{Feature engineering}: \py{pd.cut()}, \py{pd.qcut()}, column combination.
    \item \textbf{Encoding}: \py{pd.get\_dummies()} (one-hot), \py{LabelEncoder} (ordinal).
    \item \textbf{Scaling}: $z = \frac{x - \mu}{\sigma}$ (StandardScaler), $x' = \frac{x - x_{\min}}{x_{\max} - x_{\min}}$ (MinMaxScaler).
    \item \textbf{Golden rule}: document every transformation and avoid \emph{data leakage}.
\end{itemize}
\end{keyformulas}
